% Wersja źródłowa: audyt II, jednostka 30
\setcounter{section}{29}
\noindent\textbf{Wyniki wykorzystywane.} Przyjmujemy znane twierdzenie Hahna--Banacha,
lemat Riesza, twierdzenie o odwzorowaniu otwartym, własności operatorów
zwartych i twierdzenie Schaudera o zwartości transpozycji.
\par\medskip

\section{Teoria Fredholma i twierdzenie Riesza--Schaudera}
\label{sec:fredholm30}

W tym rozdziale $X,Y$ są przestrzeniami Banacha nad
$\K\in\{\R,\C\}$, a $\mathcal L(X,Y)$ i $\mathcal K(X,Y)$
oznaczają odpowiednio operatory liniowe ciągłe i zwarte.
Kryteria zwartości i twierdzenie Schaudera o zwartości transpozycji
zostały już udowodnione. Najpierw przedstawimy alternatywę
Fredholma w przestrzeni Hilberta, w wersji korzystającej z
przybliżeń przez rzuty ortogonalne, a następnie argument
działający w dowolnych przestrzeniach Banacha.

\subsection{Operatory Fredholma i ich indeks}

\begin{definicja}[Operator Fredholma]
\label{def:fredholm30}
Niech $X,Y$ będą przestrzeniami Banacha i $T\in\mathcal L(X,Y)$.
Operator $T$ nazywamy \emph{operatorem Fredholma}, jeżeli
$\operatorname{ran}T$ jest domknięty, $\dim\ker T<\infty$ i
$\operatorname{codim}\operatorname{ran}T:=\dim(Y/\operatorname{ran}T)<\infty$.
Jego \emph{indeks Fredholma} wynosi
\[\operatorname{ind}T=\dim\ker T-\operatorname{codim}\operatorname{ran}T.\]
Symbol $X'$ oznacza przestrzeń dualną topologiczną $X$.
\end{definicja}

\begin{lemat}[Uzupełniające podprzestrzenie]
\label{lem:fredholm-complements30}
Każda skończenie wymiarowa podprzestrzeń $N$ przestrzeni Banacha
$X$ jest domknięta i ma domknięte dopełnienie topologiczne
$X=N\oplus X_1$, z ciągłymi rzutami. Jeżeli $R\subseteq Y$
jest domkniętą podprzestrzenią Banacha o skończonym kowymiarze,
to $Y=R\oplus C$, gdzie $C$ ma skończony wymiar i oba rzuty
są ciągłe.
\end{lemat}
\begin{proof}
Dla bazy $u_1,\ldots,u_n$ przestrzeni $N$ funkcjonały
współrzędnych są ciągłe na $N$ i przedłużają się na $X$
przez twierdzenie Hahna--Banacha do $f_1,\ldots,f_n\in X'$.
Wówczas $P(x)=\sum_{j=1}^n f_j(x)u_j$ jest ciągłym rzutem
na $N$, a $X_1=\ker P$. W drugim przypadku wybieramy bazę
$[v_1],\ldots,[v_m]$ skończenie wymiarowego ilorazu $Y/R$
i definiujemy $C=\operatorname{span}\{v_1,\ldots,v_m\}$.
Odwzorowanie przypisujące $y\in Y$ jego składową w $C$
przez ciągłe funkcjonały współrzędnych $Y/R$ jest ciągłym
rzutem o jądrze $R$. Obie dekompozycje są topologiczne.
\end{proof}

\begin{twierdzenie}[Obraz domknięty i warunek zgodności]
\label{thm:fredholm-compatibility30}
Jeżeli $T\in\mathcal L(X,Y)$ ma domknięty obraz, to
\[\operatorname{ran}T=\{y\in Y:\ f(y)=0\ \text{dla każdego }f\in\ker T'\}.\]
Jeżeli ponadto $T$ jest operatorem Fredholma, to
$\dim\ker T'=\operatorname{codim}\operatorname{ran}T$.
Zatem $Tx=y$ ma rozwiązanie wtedy i tylko wtedy, gdy
$f(y)=0$ dla wszystkich $f\in\ker T'$, a zbiór
rozwiązań ma postać $x_0+\ker T$.
\end{twierdzenie}
\begin{proof}
Z definicji $T'$ mamy $\ker T'=(\operatorname{ran}T)^\perp$,
gdzie symbol $R^\perp=\{f\in Y':f|_R=0\}$ oznacza
anihilator. Gdy $R=\operatorname{ran}T$ jest domknięty,
twierdzenie Hahna--Banacha oddziela każdy $y\notin R$
od $R$ ciągłym funkcjonałem zerującym się na $R$.
Stąd $R=(R^\perp)_\perp$. Ponadto $R^\perp$ jest
izomorficzna z $(Y/R)'$, które przy skończonym wymiarze
$Y/R$ ma ten sam wymiar. Reszta wynika z liniowości.
\end{proof}

\subsection{Alternatywa Fredholma w przestrzeni Hilberta}

\begin{twierdzenie}[Alternatywa Fredholma dla operatora zwartego]
\label{thm:fredholm-hilbert30}
Niech $H$ będzie zespoloną przestrzenią Hilberta, a
$K\in\mathcal K(H,H)$. Wówczas $T=I_H-K$ ma domknięty
obraz oraz skończenie wymiarowe jądra $\ker T$ i $\ker T^*$,
przy czym $\dim\ker T=\dim\ker T^*$.
Dla $f\in H$ równanie $Tu=f$ ma rozwiązanie dokładnie wtedy,
gdy $f\perp\ker T^*$; wszystkie rozwiązania tworzą
klasę $u_0+\ker T$.
\end{twierdzenie}
\begin{proof}
\emph{Krok 1: skończony wymiar jądra.}
Gdyby $N=\ker T$ był nieskończenie wymiarowy, zawierałby
nieskończony układ ortonormalny $(e_n)$. Równość $Ke_n=e_n$
przeczyłaby zwartości $K$.

\emph{Krok 2: domkniętość obrazu.}
Niech $M=N^\perp$. Istnieje $c>0$ takie, że
$\|Tx\|\ge c\|x\|$ dla $x\in M$: w przeciwnym przypadku
istniałby ciąg $x_n\in M$ o normach $1$, dla którego
$Tx_n\to0$. Ze zwartości $K$ wybralibyśmy podciąg
$(x_{n_j})$ z $Kx_{n_j}\to z$, a wtedy
$x_{n_j}=Tx_{n_j}+Kx_{n_j}\to z\in M\cap N$, przy $\|z\|=1$.
To sprzeczność. Jeżeli $Tx_n\to y$, zastępujemy $x_n$
jego rzutem ortogonalnym na $M$; oszacowanie dolne
czyni ten ciąg ciągiem Cauchy'ego. Jego granica $x\in M$
spełnia $Tx=y$.

\emph{Krok 3: indeks zerowy.}
Zwartość $K(B_H)$ daje skończenie wymiarową podprzestrzeń
$W\subseteq H$, dla której
$\|K-P_WK\|<1$: wystarczy wziąć podprzestrzeń
rozpiętą na punktach odpowiednio dokładnej skończonej siatki.
Piszemy $F=P_WK$, $E=K-F$, $\|E\|<1$ i
$G=(I-E)^{-1}F$. Szereg Neumanna dowodzi, że $I-E$
jest odwracalny, a
\[I-K=(I-E)(I-G).\]
Operator $G$ ma skończony rząd. Niech $W'=\operatorname{ran}G$
i rozłóżmy $H=W'\oplus(W')^\perp$.
W tych współrzędnych
\[(I-G)(w,z)=((I_{W'}-G|_{W'})w-Gz,\ z).\]
Składowa druga identyczności pozwala eliminować $-Gz$
i redukuje obliczenie wymiaru jądra i kowymiaru obrazu
do skończenie wymiarowego operatora $I_{W'}-G|_{W'}$.
Na mocy twierdzenia o wymiarze jądra i obrazu obie liczby są
równe. Mnożenie z lewej przez $I-E$ nie zmienia
indeksu, więc $\operatorname{ind}T=0$.

\emph{Krok 4: rozwiązywalność.}
Ponieważ $\operatorname{ran}T$ jest domknięty,
$(\operatorname{ran}T)^\perp=\ker T^*$ daje
$\operatorname{ran}T=(\ker T^*)^\perp$. Z kroku 3 wynika
$\dim\ker T=\dim\ker T^*$, zaś opis rozwiązań
wynika z liniowości.
\end{proof}

\begin{uwaga}[Zakres powyższego dowodu]
\label{rem:fredholm-hilbert-scope30}
W dowodzie wykorzystano rzut ortogonalny na skończenie
wymiarową podprzestrzeń. Dowolna przestrzeń Banacha
nie musi mieć własności aproksymacji, dlatego nie wolno
przenosić tego argumentu bez zmian. Alternatywę Fredholma
w przestrzeniach Banacha udowodnimy bez takiej aproksymacji.
\end{uwaga}

\subsection{Twierdzenie Atkinsona i indeks w przestrzeniach Banacha}

\begin{twierdzenie}[Twierdzenie Atkinsona]
\label{thm:atkinson30}
Dla $T\in\mathcal L(X,Y)$ następujące warunki są równoważne:
\begin{enumerate}[label=\textup{(\roman*)}]
\item $T$ jest operatorem Fredholma;
\item istnieje $S\in\mathcal L(Y,X)$, dla którego
$ST-I_X\in\mathcal K(X,X)$ i $TS-I_Y\in\mathcal K(Y,Y)$.
\end{enumerate}
Gdy (i) zachodzi, obie różnice można wybrać skończonego rzędu.
\end{twierdzenie}
\begin{proof}
Załóżmy (i) i wybierzmy rozkłady $X=N\oplus X_1$,
$Y=R\oplus C$, gdzie $N=\ker T$, $R=\operatorname{ran}T$.
Ograniczenie $A=T|_{X_1}:X_1\to R$ jest bijekcją między
przestrzeniami Banacha, więc $A^{-1}$ jest ograniczony.
Dla $y=r+c$ definiujemy $S(y)=A^{-1}r$.
Wtedy $ST=I_X-P_N$ i $TS=I_Y-P_C$, gdzie oba rzuty
mają skończony rząd.

Odwrotnie niech $ST=I_X-K_1$ i $TS=I_Y-K_2$,
gdzie $K_1,K_2$ są zwarte. Na $N=\ker T$ operator
$K_1$ działa jak identyczność, więc $N$ jest skończenie
wymiarowe (inaczej lemat Riesza przeczyłby zwartości).
Wybierzmy dopełnienie $X=N\oplus X_1$.
Jeśli $T|_{X_1}$ nie byłby ograniczony z dołu, istniałby
ciąg jednostkowy $x_n\in X_1$ z $Tx_n\to0$.
Równość $x_n=STx_n+K_1x_n$ oraz zwartość $K_1$
dawałyby normowo zbieżny podciąg do niezerowego wektora
w $N\cap X_1$, co jest niemożliwe. Zatem obraz $T$
jest domknięty.
Z równości $S'T'=I_{Y'}-K_2'$ (otrzymanej przez
transpozycję $TS=I_Y-K_2$) wynika, że
$\ker T'\subseteq\ker(I_{Y'}-K_2')$.
Operator $K_2'$ jest zwarty na mocy twierdzenia Schaudera,
więc drugie jądro ma skończony wymiar.
Twierdzenie~\ref{thm:fredholm-compatibility30} daje
$\operatorname{codim}\operatorname{ran}T=\dim\ker T'<\infty$.
\end{proof}

\begin{twierdzenie}[Lokalna stałość indeksu]
\label{thm:index-local30}
Niech $T\in\mathcal L(X,Y)$ będzie operatorem Fredholma.
Istnieje $\delta>0$ takie, że każdy $U\in\mathcal L(X,Y)$
z $\|U-T\|<\delta$ jest Fredholma i
$\operatorname{ind}U=\operatorname{ind}T$.
\end{twierdzenie}
\begin{proof}
Wybierzmy rozkłady $X=N\oplus X_1$, $Y=R\oplus C$
z dowodu twierdzenia Atkinsona. W ich współrzędnych
$T(n,x)=(Ax,0)$, gdzie $A:X_1\to R$ jest izomorfizmem.
Dla $U$ bliskiego $T$ mamy
\[U(n,x)=(Bn+Dx,\ En+Fx),\]
gdzie $B:N\to R$, $D:X_1\to R$, $E:N\to C$,
$F:X_1\to C$. Operator $D$ jest odwracalny, gdy
$\|A^{-1}(D-A)\|<1$, co zachodzi dla $U$ dostatecznie
bliskiego $T$.
Dokonujemy dwóch \emph{odwracalnych} zmian współrzędnych.
Najpierw w dziedzinie podstawiamy
$(n,x)=(n,x'-D^{-1}Bn)$. Otrzymujemy
\[U(n,x)=(Dx',\ (E-FD^{-1}B)n+Fx').\]
Następnie w przeciwdziedzinie stosujemy
$(r,c)\mapsto(r,c-FD^{-1}r)$. Wynik ma postać
\[(n,x')\mapsto(Dx',\ (E-FD^{-1}B)n).\]
Obraz jest domknięty, ponieważ $D$ jest izomorfizmem,
a $E-FD^{-1}B:N\to C$ działa między przestrzeniami
skończenie wymiarowymi. Indeks drugiego bloku to
$\dim N-\dim C$ niezależnie od jego rzędu, gdyż
\[\dim\ker M-\dim(C/\operatorname{ran}M)=\dim N-\dim C\]
dla każdego $M:N\to C$ w skończonym wymiarze.
Stąd $\operatorname{ind}U=\dim N-\dim C=\operatorname{ind}T$.
\end{proof}

\begin{twierdzenie}[Zaburzenia zwarte zachowują indeks]
\label{thm:index-compact-perturb30}
Niech $T\in\mathcal L(X,Y)$ będzie operatorem Fredholma,
a $K\in\mathcal K(X,Y)$. Wówczas $T+K$ jest Fredholma i
$\operatorname{ind}(T+K)=\operatorname{ind}T$.
\end{twierdzenie}
\begin{proof}
Z twierdzenia Atkinsona istnieje $S:Y\to X$ takie, że
$ST-I_X$ i $TS-I_Y$ są zwarte. Dla każdego $t\in[0,1]$
obie różnice $S(T+tK)-I_X$ i $(T+tK)S-I_Y$ są zwarte.
Zatem $T+tK$ jest Fredholma dla wszystkich $t$.
Jego indeks zależy lokalnie stale od $t$ na mocy
Twierdzenia~\ref{thm:index-local30}; ponieważ $[0,1]$
jest spójny, jest stały.
\end{proof}

\begin{wniosek}[Alternatywa Fredholma w przestrzeni Banacha]
\label{cor:fredholm-banach30}
Niech $X$ będzie przestrzenią Banacha i
$K\in\mathcal K(X,X)$. Wówczas $I_X-K$ jest operatorem
Fredholma indeksu zero. Równanie $(I_X-K)x=y$ jest
rozwiązywalne dokładnie wtedy, gdy $f(y)=0$ dla każdego
$f\in\ker(I_{X'}-K')$. Jeśli $\ker(I_X-K)=\{0\}$,
to $I_X-K$ jest izomorfizmem liniowo-topologicznym.
\end{wniosek}
\begin{proof}
Operator $I_X$ jest Fredholma indeksu zero; zaburzenie
$-K$ zachowuje indeks. Warunek zgodności wynika z
Twierdzenia~\ref{thm:fredholm-compatibility30}. Gdy jądro
jest zerowe, indeks zero i domkniętość obrazu dają
surjektywność. Ciągłość odwrotności wynika z twierdzenia
o odwzorowaniu otwartym.
\end{proof}

\begin{przyklad}[Indeks może być niezerowy]
\label{ex:shift-index30}
W $X=\ell^2(\N)$ przesunięcie w prawo
$R(x_1,x_2,\ldots)=(0,x_1,x_2,\ldots)$ ma zerowe jądro
i obraz kowymiaru jeden, a więc indeks $-1$.
Przesunięcie w lewo $L(x_1,x_2,\ldots)=(x_2,x_3,\ldots)$
ma indeks $1$. Nie każdy operator Fredholma ma więc
indeks zero.
\end{przyklad}

\subsection{Twierdzenie Riesza--Schaudera i widmo operatora zwartego}

\begin{definicja}[Widmo i rezolwenta]
\label{def:spectrum30}
Niech $X$ będzie zespoloną przestrzenią Banacha i
$K\in\mathcal L(X,X)$. Definiujemy
$\rho(K)=\{\lambda\in\C:\lambda I_X-K\text{ jest odwracalny}\}$
oraz $\sigma(K)=\C\setminus\rho(K)$.
Liczbę $\lambda$ nazywamy wartością własną, gdy
$\ker(\lambda I_X-K)\ne\{0\}$.
\end{definicja}

\begin{twierdzenie}[Riesz--Schauder]
\label{thm:riesz-schauder30}
Niech $X$ będzie nieskończenie wymiarową zespoloną
przestrzenią Banacha, a $K\in\mathcal K(X,X)$.
Wtedy:
\begin{enumerate}[label=\textup{(\roman*)}]
\item $0\in\sigma(K)$, a $\sigma(K)$ jest niepustym
zwartym podzbiorem $\C$;
\item każdy $\lambda\in\sigma(K)\setminus\{0\}$ jest
wartością własną o skończenie wymiarowym jądrze;
\item dla każdego $\varepsilon>0$ zbiór
$\{\lambda\in\sigma(K):|\lambda|\ge\varepsilon\}$
jest skończony;
\item dla $\lambda\ne0$ uogólniona przestrzeń własna
$\bigcup_{m\ge1}\ker(\lambda I_X-K)^m$ jest
skończenie wymiarowa i równa
$\ker(\lambda I_X-K)^{m_0}$ dla pewnego $m_0\ge1$.
\end{enumerate}
\end{twierdzenie}
\begin{proof}
Dla $|\lambda|>\|K\|$ szereg Neumanna dowodzi
odwracalności $\lambda I_X-K$. Zbiór rezolwenty jest otwarty:
jeśli $\lambda_0 I_X-K$ ma odwrotność, to dla małego
$|\lambda-\lambda_0|$ stosujemy szereg Neumanna do
$I+(\lambda-\lambda_0)(\lambda_0 I_X-K)^{-1}$.
Widmo jest więc domknięte i ograniczone. Ponieważ
$K$ nie może być odwracalny (inaczej
$I_X=K^{-1}K$ byłby zwarty), $0\in\sigma(K)$.

Dla $\lambda\ne0$ operator $\lambda I_X-K$ jest Fredholma
indeksu zero jako niezerowa wielokrotność $I_X-\lambda^{-1}K$.
Jeżeli nie jest odwracalny, nie może być różnowartościowy:
zerowe jądro i indeks zero wymuszałyby surjektywność.
Jądro jest skończenie wymiarowe. To dowodzi (ii).

Gdyby (iii) zawiodło, wybralibyśmy parami różne
niezerowe wartości własne $\lambda_n$ o
$|\lambda_n|\ge\varepsilon$ i odpowiadające im wektory
własne $v_n\ne0$. Wektory te są liniowo niezależne.
Niech $M_n=\operatorname{span}\{v_1,\ldots,v_n\}$.
Podprzestrzenie $M_n$ są niezmiennicze względem $K$.
Lemat Riesza dostarcza $x_n\in M_n$ o normie $1$ i
$\operatorname{dist}(x_n,M_{n-1})>1/2$ dla $n\ge2$.
Na ilorazie $M_n/M_{n-1}$ operator $K$ działa przez
mnożenie przez $\lambda_n$, więc
$Kx_n-\lambda_nx_n\in M_{n-1}$. Jeśli $m<n$,
to $Kx_m\in M_{n-1}$, a zatem
\[\|Kx_n-Kx_m\|\ge
|\lambda_n|\operatorname{dist}(x_n,M_{n-1})>\varepsilon/2.\]
To przeczy zwartości $K$.

Dla (iv) połóżmy $N_m=\ker(\lambda I_X-K)^m$.
Wielomian $\lambda^{-m}(\lambda I_X-K)^m$ ma postać
$I_X-Q_m$ z $Q_m$ zwartym (każdy niezerowy składnik
w $Q_m$ zawiera dodatnią potęgę $K$). Na mocy
alternatywy Fredholma każde $N_m$ ma skończony wymiar.
Gdyby rosnący ciąg $N_m$ nie stabilizował się, można
byłoby wybrać rosnące indeksy $n_j$ z
$N_{n_j-1}\subsetneq N_{n_j}$ oraz wektory
$x_j\in N_{n_j}$ o normie $1$ takie, że
$\operatorname{dist}(x_j,N_{n_j-1})>1/2$.
Dla $i<j$ zachodzą $Kx_i\in N_{n_j-1}$ i
$Kx_j-\lambda x_j\in N_{n_j-1}$; stąd
$\|Kx_j-Kx_i\|>|\lambda|/2$, co przeczy zwartości.
Zatem ciąg jąder stabilizuje się.
\end{proof}

\begin{przyklad}[Widmo operatora diagonalnego]
\label{ex:diagonal-spectrum30}
W $\ell^2(\N)$ operator
$K(x_1,x_2,\ldots)=(x_1,x_2/2,x_3/3,\ldots)$
jest zwarty jako granica operatorów skończonego rzędu.
Ma widmo $\sigma(K)=\{0\}\cup\{1/n:n\ge1\}$;
każda niezerowa wartość własna jest prosta.
\end{przyklad}

\begin{przyklad}[Operator Volterry: widmo bez niezerowych wartości własnych]
\label{ex:volterra-spectrum30}
W $X=C([0,1],\C)$ definiujemy
$(Vf)(t)=\int_{[0,t]}f(s)\,\mathrm d\lambda_1(s)$.
Operator $V$ jest zwarty. Iterując wzór na całkę,
dostajemy $\|V^n\|\le 1/n!$, ponieważ jądro $V^n$
jest równe $(t-s)^{n-1}/(n-1)!$ dla $0\le s\le t$.
Dla $\lambda\ne0$ szereg
$\lambda^{-1}\sum_{n\ge0}(V/\lambda)^n$
jest zbieżny w normie operatorowej i określa
$(\lambda I_X-V)^{-1}$. Zatem $\sigma(V)=\{0\}$,
mimo że $V\ne0$.
\end{przyklad}

\subsection{Zastosowanie w teorii równań różniczkowych cząstkowych}

\begin{uwaga}[Od operatora eliptycznego do operatora zwartego]
\label{rem:elliptic-fredholm30}
W skrypcie \emph{Równania różniczkowe cząstkowe, tom I}
rozpatruje się faktoryzację $L=L_\mu(I_V-K)$,
w której $L_\mu:V\to V^*$ jest izomorfizmem,
a $K:V\to V$ operatorem zwartym otrzymanym dzięki
zwartemu zanurzeniu Sobolewa. Wówczas $\operatorname{ind}L=0$,
a warunek rozwiązywalności wynika z anihilatora obrazu.
Wyprowadzenie odpowiednich nierówności koercywności
i zagadnienia brzegowe pozostają w skrypcie PDE.
\end{uwaga}
